\chapter{Conclusion}
\label{ch:conclusion}

% 독립 연구의 결론을 반복하지 않고, 두 연구축이 하나의 Claim-Aligned
% Assurance 논문을 이루는 이유와 각 보증의 중단 경계를 중심으로 정리합니다.
\section{Summary}
\label{sec:ch5-summary}

Trust claims about artificial intelligence systems are often communicated through
compact labels, parameters, or acceptance decisions.

These outputs can hide the premises that give a claim its meaning.  A differential privacy parameter is
incomplete without its protected unit, adjacency, transformation, mechanism, and
accountant; a watermark score is not by itself evidence that embedding occurred
under a declared generation context.  This dissertation addressed the gap between
such public statements and the mechanisms and evidence that warrant them.

The resulting perspective is \ClaimAlignedAssurance{}.  A positive statement is
released only when its target, operative mechanism, verification interface, and
evidence refer to the same protected or verified object.  An incomplete chain
yields a narrower qualified statement or an explicit block.  The framework is
instantiated in two technically distinct settings: evidence-bound auditing of a
completed privacy execution and proof-bound public verification of generative-image
provenance.

Chapter~\ref{ch:post-execution-validation} treated a reported privacy sentence as
a query over frozen evidence rather than as a synonym for an accountant output.
The validator reopens the raw-to-generated mapping, checks pipeline and runtime
relations, recomputes the registered accountant, resolves the requested unit and
adjacency, and releases only the exact native or converted wording licensed by the
result.  The complete contract returned the expected tuple for all 100 allowed and
100 paired non-allow cases.  Fixed necessity witnesses and preserved development
failures exposed why mapping, stability, runtime, vacuity, and exact wording cannot
be collapsed into one scalar check.  WISDM and UCI HAR executions demonstrated
query-dependent authorization, while Sepsis and Backblaze rescans showed that
window policy can substantially change raw-to-generated influence without itself
constituting a DP certificate.

Chapter~\ref{ch:pp-mark} changed the assurance object to generative image
provenance.  PP-Mark binds a latent watermark to generation context, commits a
sampled embedding trace, and verifies an image-bound opening through an SP1
receipt.  PP-Mark(score) provides transformation tolerant statistical screening;
PP-Mark(accept) requires both the calibrated score and a valid receipt for the same
context and exact image.  A 10,000-binding search and white-box optimization
exposed score-only acceptance, whereas the evaluated forgery conditions produced
no full acceptances.  Score-mode true positive rates were 0.74--1.00 over the
reported individual transformations and 0.98 after regeneration.  These are
observed results under the evaluated protocols, not a universal zero forgery
claim.

The two studies are not one algorithm.  The privacy audit limits what may be said
about an existing execution, whereas PP-Mark changes the evidence required by a
public provenance predicate.  Their common result is that a defensible claim must
name its object, bind the operative mechanism, expose an interface with fixed
semantics, and carry the evidence required by those semantics.


\section{Integrated Discussion and Contributions}
\label{sec:ch5-integrated-discussion}

\subsection{Auditable Privacy: Evidence-Relative Post-Execution Audit}
\label{subsec:ch5-auditable-privacy}

The privacy contribution begins from a practical condition: training has already
finished and the available evidence is fixed.  Different stakeholders may still
ask whether the same model supports a generated-window, raw-event, patient, owner,
or device-level sentence.  Those requests are different mathematical claims even
when they display the same base epsilon.  The query must therefore be a scientific
input to the audit rather than a label selected after one evidence-only verdict.

The validator connects the requested sentence through a chain of raw adjacency,
stable transformation, executed mechanism, accountant, and exact wording.  It
separates the derivation path, public release state, assurance authority, numerical
consequence, and issue set.  This distinction permits one frozen execution to
license a native claim, license a weaker converted claim, and block a stronger
claim.  A blocked result does not prove that the stronger privacy property is
false; it states that the supplied evidence and registered procedure do not
license that sentence.

The experimental contribution is correspondingly about decision validity and
authority.  The 200-case oracle tests controlled conformance rather than estimating
error prevalence.  WISDM supplies a complete three-query execution; UCI HAR shows
that missing raw support and owner bounds correctly stop the conclusion at the
published-window unit.  Sepsis and Backblaze expose mapping differences at medical
and operational scale, while intentionally withholding privacy wording because no
executed mechanism and accountant are attached.  Separate verifiers establish
software and artifact separation from the production path, not execution
attestation.

This produces a concrete reuse boundary.  If an existing model's frozen evidence
licenses the requested unit, the model can be reused with the exact supported
wording.  If a registered conversion produces a nonvacuous pair, only that weaker
pair may be stated.  If the requested patient or owner claim is blocked, the base
window pair cannot be relabeled; a mechanism and execution native to the stronger
unit may be required.  Determining the cost and utility of that alternative is a
future experiment, not a result silently inferred from the audit.

\subsection{Secure Public Verification}
\label{subsec:ch5-secure-public-verification}

The provenance study exposed a tension between independent verification and
attack surface.  A private detector requires trust in its operator; a public
detector enables independent screening but exposes a queryable or differentiable
decision surface.  PP-Mark therefore does not claim to make the score invulnerable.
It changes the stronger acceptance rule so that optimizing the public statistic is
insufficient without a valid context-, image-, and trace-bound receipt.

The score and receipt establish different facts.  The score supplies image-side
evidence that can survive selected lossy transformations.  The receipt supplies
process-side evidence for the registered binding, trace commitment, sampled
embedding relation, and exact-image predicate.  Full acceptance requires both,
allowing third-party verification without a private API or disclosure of the
producer secret.  This construction is not a replacement for a conventional
signature when exact integrity alone is required; its added value is verification
of the embedding relation together with a durable in-image signal.

The interface deliberately has three outcomes rather than an implicit fallback.
An exact image with a valid receipt can receive PP-Mark(accept).  A transformed
copy may retain PP-Mark(score) but does not inherit exact-image acceptance.  If the
public binding or receipt cannot be resolved, the result is unverified rather than
negative evidence that an image is not AI-generated.  PP-Mark is therefore a
positive provenance system for participating pipelines, not a universal synthetic-
image classifier.

\subsection{Claim-Aligned Assurance Design Principles}
\label{subsec:ch5-caa-principles}

The integrated analysis supports four dissertation-wide principles.

\begin{enumerate}
    \item \textbf{Claim specification.}  A positive statement identifies its
    target and semantics: the privacy unit, adjacency, and wording, or the exact
    image, generation context, and verification mode.

    \item \textbf{Mechanism binding.}  The statement is connected to the mechanism
    that gives it meaning.  Privacy binds transformation, sampling, clipping,
    noise, accounting, and runtime evidence; provenance binds payload, latent
    embedding, trace commitment, image-bound challenge, and proof predicate.

    \item \textbf{Evidence-bound decision.}  An output does not inherit the
    authority of missing evidence.  Positive privacy wording requires its mapping,
    runtime, and accounting relations; full provenance acceptance requires both
    image-side detection and process-side receipt verification.

    \item \textbf{Fail-closed assurance.}  A missing, inconsistent, invalid,
    unregistered, or vacuous premise blocks the stronger interpretation.  Blocking
    does not prove the broader property false; it states that the requested
    conclusion is not licensed by the declared procedure.
\end{enumerate}

The verification interface operationalizes all four principles.  Query separation
prevents a window result from becoming a patient answer, while separate score and
acceptance modes prevent statistical detection from becoming process-bound
provenance.  The reusable design question is consequently narrow: which observed
or proven relations license the exact positive sentence being released?


\section{Limitations, Broader Impact, and Future Work}
\label{sec:ch5-limits-impact-future}

\subsection{Limitations and Assurance Boundaries}
\label{subsec:ch5-limitations}

The dissertation covers two assurance domains and two principal systems.  It does
not establish a complete framework for all trustworthy-AI claims.  Fairness,
safety, factuality, utility, and governance can require different objects and
evidence.  The four principles are a design abstraction grounded in the evaluated
settings, not a universal certification standard.

The privacy result remains evidence-relative and finite.  It assumes a fallible
but nonmalicious producer, registered transformations and accountant adapters, and
faithful disclosure of the observable bundle.  A malicious producer can fabricate
a mutually consistent bundle after executing another procedure.  The validator
cannot distinguish those worlds without authenticated observation, and therefore
does not attest execution, prove arbitrary program semantics, or replace the
underlying DP proof.  The controlled oracle establishes exact behavior on its
authored 200 cases but does not estimate how often published systems contain such
errors.  Sepsis and Backblaze are mapping diagnostics only, not trained DP models
or patient-level certificates.

PP-Mark is conditional on hash resistance, proof-system soundness, producer-key
control, inversion behavior, the composite predicate, and the partial-opening
bound.  It is evaluated on Stable Diffusion 2.1 and SDXL, two latent-diffusion
models that share a compatible embedding and inversion interface.  The submitted
SP1 v5 implementation requires about 48.6 seconds and produces a 40.68 MB receipt
at 1,000 sampled positions.  A post-submission SP1 v6.3.1 port reduced the measured
values to 17.0 seconds and 2.81 MB, but this update does not eliminate proof cost or
establish a universal deployment profile.  Exact-image acceptance is invalidated
by pixel edits, the public binding and receipt must remain retrievable, and key
rotation, revocation, geometric synchronization, and broader model families remain
open.

Claim alignment can prevent a system from saying more than its evidence supports,
but it cannot repair false source evidence, a compromised key, or facts outside a
proof circuit.  It narrows the statement and makes the stopping boundary visible;
it does not eliminate every trust assumption.

\subsection{Broader Impact}
\label{subsec:ch5-broader-impact}

Auditable Privacy applies to sequential learning with medical sensor data,
clinical records, activity data, and operational logs.  Several privacy units can
be plausible in each setting.

Explicit units, transformations, runtime relations, and evidence can
help developers and auditors avoid presenting generated record accounting as a
stronger patient-, user-, or device-level promise.  It can also clarify whether an
existing model is reusable under narrower wording or whether a new native unit
execution is required.  This remains an audit of claim support, not clinical
certification.

Secure Public Verification can support creators, platforms, archives, and
reviewers that verify declared provenance without one private verification API.
Its separated modes clarify that a transformed image can retain a detectable
signal while requiring a new receipt for exact acceptance.  Provenance evidence
nevertheless does not prove factual accuracy, human authorship, social legitimacy,
or legal ownership.

Fail-closed systems also impose costs.  Unavailable evidence may withhold a
legitimate claim, proof expense may exclude producers, and key ownership may
centralize control.  Deployment therefore requires accessible verification,
transparent key governance, clear communication of unverified outcomes, and an
appeal or re-issuance process.

\subsection{Future Experimental Validation}
\label{subsec:ch5-future-validation}

The first future study will quantify the boundary between reuse through post-hoc
conversion and retraining with native patient-level DP.  The study will use the
5,000-patient PhysioNet/CinC 2019 Sepsis subset already audited in
Chapter~\ref{ch:post-execution-validation}, containing 370 sepsis-positive and
4,630 control patients.  A source- and label-stratified patient split will allocate
3,500 patients to training (259 positive, 3,241 control), 750 to validation (55
positive, 695 control), and 750 to testing (56 positive, 694 control), with no
patient crossing a split.

Three training conditions will be compared on the same patient partitions: a
nonprivate reference, window-level DP-SGD followed by the evidence-bound audit,
and native patient-level DP-SGD.  The native condition will sample patients,
aggregate each selected patient's window gradients, clip once per patient, add
Gaussian noise, and account at patient adjacency.  Target patient-level epsilon
values of 2, 4, and 8, public window-contribution caps of 1, 4, 8, and 16, and five
fixed seeds will define the comparison grid.  Evaluation will report AUPRC,
calibration and utility, the licensed or blocked post-hoc privacy statement,
membership-inference exposure, training time, and peak memory.  The intended output
is a measured rule for when an existing window-private model remains reusable and
when a native patient-private execution is necessary.

The second future study will make PP-Mark proof artifacts discoverable after
distribution.  A C2PA-compatible external manifest repository and resolver will
store or reference the public binding \(\PPBinding\), PP-Mark public statement,
SP1 receipt, producer signature, model and inversion configuration, and revocation
metadata.  C2PA is used here as a signed transport and discovery layer; the design
does not assume that C2PA provides one universal central registry or proves the
PP-Mark embedding relation.

An exact image will first use a manifest identifier, URI, or hard binding to
retrieve its proof package.  When embedded metadata has been stripped, a soft
fingerprint resolver will return candidate manifests for cryptographic and
statistical verification rather than directly declaring provenance.  Experiments
will include metadata stripping, JPEG compression, crop--resize, manifest
substitution, cross-image receipt replay, missing records, and key revocation.
Metrics will include retrieval Recall@\(k\), false-acceptance and true-positive
rates, p95 lookup and verification latency, registry storage, and stale-record
rejection.  The decision policy remains explicit: an exact image with a valid
score and proof receives \textsf{Accept}; a transformed image may receive
\textsf{Score}; and an unresolved artifact returns \textsf{Unverified}.

Both future studies close concrete gaps exposed by the completed work.  The first
tests the cost of moving from an audited converted claim to a native protected unit.
The second tests how proof-bound public verification survives real distribution
when metadata and sidecars may be removed.  Neither result is claimed as completed
in this dissertation.

Within the scope evaluated here, the dissertation reaches a narrower conclusion.
Privacy and provenance claims become defensible not when a system exposes the
largest number of metrics, but when each positive statement is limited to the
object, mechanism, interface, and evidence that jointly support it.  That is the
role of Claim-Aligned Assurance in trustworthy AI.
